\chapter{Spherical converse, stage 2: mixed-sign curvature}
\label{chap:sphere_mixed}

% Chapter for Gluck/SphereMixed.lean (Stage 2 of the S² extension).
% Design-level blocks authored iter-048 from the S2-C/S2-D interfaces of
% chap:sphere; declaration names are [expected] pins (the Lean file does not
% exist yet — tex precedes Lean). Refine after S2-D lands.
%
% SOURCE background: references/spaceform_notes.md, sections "THE CRUX AND ITS
% RESOLUTION", "H² vs S²" (read iter-048). Key asymmetry: on S² hemisphere
% confinement does the topology job but NOT the positive-speed job; the true
% admissible range allows κ_S ≤ 0 where ⟨z,n⟩ > 0, quantitatively κ_S > −R on
% the confinement disk |z| ≤ R. This chapter states the mixed-sign converse
% through that mechanism.

\section{Overview}

Stage 2 removes global positivity from \cref{thm:spherical_converse_pos}: the
prescribed geodesic curvature $\kappa$ may be $\le0$ on part of the circle,
provided the position-dependent admissibility $\kappa(\theta)-\langle
z(\theta),ie^{i\theta}\rangle>0$ can be maintained. The pivotal observation
(iter-048) is that the \emph{entire} S2-D winding layer of \cref{chap:sphere}
is sign-agnostic and is reused verbatim:

\begin{itemize}
\item The four-vertex overlap window
  $\bigl(\max(\kappa(q_1),\kappa(q_2)),\,\min(\kappa(p_1),\kappa(p_2))\bigr)$
  has a \emph{positive} upper part whenever the two maxima are positive — even
  when the minima are $\le0$. The small-contrast step levels
  $a=c-h/2$, $b=c+h/2$ of \cref{lem:exists_abab_levels} are chosen inside that
  positive part, so the step model, the arc algebra, the first-variation
  expansion (\cref{lem:step_error_expansion}) and the boundary winding are
  literally the positive-stage objects.
\item The transport lemmas (\cref{lem:invariant_admissible_arc},
  \cref{lem:step_model_transport}) already take the actual curvature only
  through a pointwise lower bound $\kappa\ge\kappa_0$ \emph{with no sign
  assumption on $\kappa_0$}: the model margin hypothesis
  $\langle z^\ast(\theta),ie^{i\theta}\rangle\le\kappa_0-\delta-\mu$ is exactly
  the sign-changing admissibility in disguise. Along the near-centered-circle
  model, $\langle z^\ast,ie^{i\theta}\rangle\approx-r^\ast$, so the margin is
  satisfiable precisely when $\kappa_0>-r^\ast$ (up to slack) — the
  quantitative version of the notes' ``$\kappa_S>-R$'' lower range.
\end{itemize}

What genuinely changes: (a) the hypothesis package
(\cref{def:mixed_sign_sphere_four_vertex}); (b) the margins lemma must admit
$\kappa_0\in(-r^\ast,c]$ instead of $(0,c]$
(\cref{lem:mixed_sphere_margin_transport}); (c) two positive-stage lemmas
carry an unused positivity hypothesis (\texttt{IsCurvatureFunction}) that must
not survive into stage 2. For \cref{lem:reconstruction_ode} the relaxation to
``continuous and $2\pi$-periodic'' is folded into the stage-1 S2-E re-sign
(positivity of the speed comes from the admissibility margin, not from
$\kappa>0$; see the statement note there) — one declaration serves both
stages, no variant block needed. For \cref{lem:step_L1_reparam} the Euclidean
plateau engine \cref{lem:exists_preliminary_reparam} is frozen with its
positive hypothesis, so stage 2 uses the constant-shift reduction
\cref{lem:step_L1_reparam_relaxed} below — a genuine (small) new lemma, not an
edit to frozen files.

\section{Declarations}

\begin{definition}[mixed-sign spherical four-vertex hypothesis]
  \label{def:mixed_sign_sphere_four_vertex}
  \lean{Gluck.MixedSignSphereFourVertex}
  \uses{def:four_vertex_condition, def:sphere_four_vertex}
  A function $\kappa:\mathbb R\to\mathbb R$ satisfies the \emph{mixed-sign
  spherical four-vertex hypothesis} if it is continuous, $2\pi$-periodic, and
  either constant positive, or has the value-separated alternating extrema of
  \cref{def:four_vertex_condition} with both maxima positive
  ($\kappa(p_1),\kappa(p_2)>0$), and moreover for some $c$ in the positive
  part of the overlap window
  ($\max(0,\kappa(q_1),\kappa(q_2))<c<\min(\kappa(p_1),\kappa(p_2))$) the
  global lower bound
  \[
    \kappa(\theta)\;>\;-\bigl(\sqrt{1+c^2}-c\bigr)\qquad\text{for all }\theta
  \]
  holds. (No global positivity: $\kappa$ may be $\le0$ at and around the
  minima. The lower bound is the confinement-forcing condition: $r^\ast(c)=
  \sqrt{1+c^2}-c$ is the model-circle radius, and $\kappa>-r^\ast$ is what
  keeps the position-dependent denominator $\kappa-\langle z,ie^{i\theta}
  \rangle$ positive along trajectories near the model. Compare
  \texttt{MixedSignFourVertex} in the Euclidean Dahlberg extension, where the
  analogous package is positivity at the two maxima only.)
\end{definition}

\begin{definition}[spherical admissible profile]
  \label{def:spherical_admissible_profile}
  \lean{Gluck.SphericalAdmissibleProfile}
  \uses{def:mixed_sign_sphere_four_vertex}
  The compact data package extracted from
  \cref{def:mixed_sign_sphere_four_vertex} that the mixed-sign argument
  threads through the flow layer: the window midpoint $c$, the disk radius
  $R<1$, the denominator margin $\delta>0$, the transport margin $\mu>0$, and
  the pointwise lower bound $\kappa_0\in(-r^\ast(c),c]$ with
  $\kappa\ge\kappa_0$. (Replaces the positive-stage shortcut
  $\kappa_0=\min\kappa>0$ of \cref{lem:sphere_curvature_lower_bound}; it is
  the stage-2 stand-in the positive-stage docstrings promised.)
\end{definition}

\begin{lemma}[$L^1$ step reparametrization without positivity]
  \label{lem:step_L1_reparam_relaxed}
  \lean{Gluck.exists_step_L1_reparam_relaxed}
  \uses{lem:step_L1_reparam, def:four_arc_step_function}
  Let $\kappa$ be continuous and $2\pi$-periodic (no positivity), $0<a<b$, and
  $\theta_1<\theta_2<\theta_3<\theta_4<\theta_1+2\pi$ with
  $\kappa(\theta_1,\theta_2,\theta_3,\theta_4)=(a,b,a,b)$. For every
  $\varepsilon>0$ there is an orientation-preserving circle reparametrization
  $h_1$ ($\mathrm{StrictMono}$, $C^1$ with continuous positive derivative,
  $h_1(\theta+2\pi)=h_1(\theta)+2\pi$) such that
  \[
    \int_0^{2\pi}\bigl|\kappa(h_1(\theta))-\kappa^\ast(\theta)\bigr|\,d\theta
    \;<\;\varepsilon,
    \qquad \kappa^\ast=\texttt{stepCurvature } b\,a\,0\,(\pi/2)\,\pi\,(3\pi/2).
  \]
\end{lemma}

\begin{proof}
  \uses{lem:step_L1_reparam}
  Constant-shift reduction to the positive-stage lemma; the frozen Euclidean
  plateau engine is applied only to a positive function. Pick
  $M>\max\bigl(0,-\min_{[0,2\pi]}\kappa\bigr)$ (finite by compactness and
  periodicity), so $\kappa+M$ is a curvature function
  (continuous, $2\pi$-periodic, $>0$). The crossing data shifts to
  $(\kappa+M)(\theta_k)=(a+M,b+M,a+M,b+M)$ with $0<a+M<b+M$, so
  \cref{lem:step_L1_reparam} applied to $\kappa+M$ with levels $a+M,b+M$
  yields $h_1$ with
  $\int_0^{2\pi}|(\kappa+M)\circ h_1-\kappa^\ast_M|<\varepsilon$, where
  $\kappa^\ast_M=\texttt{stepCurvature}\,(b+M)\,(a+M)\,0\,(\pi/2)\,\pi\,
  (3\pi/2)$. Pointwise $\kappa^\ast_M=\kappa^\ast+M$ (the step function takes
  only the two level values, each shifted by $M$), so the integrand equals
  $|\kappa\circ h_1-\kappa^\ast|$ pointwise and the same $h_1$ witnesses the
  claim. ($\sim$30--60 LOC; the only Lean content beyond the application is
  the pointwise shift identity for \texttt{stepCurvature}.)
\end{proof}

\begin{lemma}[mixed-sign margin transport]
  \label{lem:mixed_sphere_margin_transport}
  \lean{Gluck.mixedSphere_margin_transport}
  \uses{def:spherical_admissible_profile, lem:invariant_admissible_arc,
        lem:step_model_transport, lem:step_model_margins}
  \Cref{lem:step_model_margins} holds with the hypothesis $\kappa_0\in(0,c]$
  relaxed to $\kappa_0\in(-r^\ast(c)+\sigma,\,c]$ for any slack $\sigma>0$
  (with margins now depending on $\sigma$): explicit $R<1$, $\delta>0$,
  $\mu>0$, $\rho_0>0$, $h_0>0$ such that the four quarter arcs of the
  concatenated step model near the centered circle satisfy
  $\|z^\ast\|\le R-\mu$,
  $\langle z^\ast(\theta),ie^{i\theta}\rangle\le\kappa_0-\delta-\mu$, and
  clamp-inactivity. Consequently \cref{lem:step_model_transport} applies to
  mixed-sign $\kappa\ge\kappa_0$ unchanged.
\end{lemma}

\begin{proof}
  Along the centered circle $\langle z^\ast(\theta),ie^{i\theta}\rangle\equiv
  -r^\ast$, so the required margin inequality reads
  $-r^\ast\le\kappa_0-\delta-\mu$, i.e.\ $\delta+\mu\le\kappa_0+r^\ast$, which
  has room exactly when $\kappa_0>-r^\ast$; take $\delta+\mu=\sigma/2$, and
  shrink $\rho_0,h_0$ by the same Lipschitz-deviation argument as in
  \cref{lem:step_model_margins}. The transport
  (\cref{lem:invariant_admissible_arc}, \cref{lem:step_model_transport}) never
  used $\kappa_0>0$.
\end{proof}

\begin{lemma}[mixed-sign endpoint winding]
  \label{lem:mixed_spherical_endpoint_winding}
  \lean{Gluck.mixed_spherical_endpoint_winding}
  \uses{def:spherical_endpoint, lem:spherical_endpoint_continuousOn,
        def:step_error_map, lem:step_model_transport,
        lem:mixed_sphere_margin_transport, def:spherical_admissible_profile,
        lem:step_error_expansion, lem:exists_abab_levels,
        lem:step_L1_reparam_relaxed, thm:existence_of_zero,
        lem:winding_conj_loop, lem:winding_number_c_perturb,
        def:mixed_sign_sphere_four_vertex}
  Let $\kappa$ satisfy \cref{def:mixed_sign_sphere_four_vertex} and be
  non-constant. Then there exist $R<1$, $\delta>0$, a circle reparametrization
  $h_1$ and $z_0$ such that the $(\kappa\circ h_1,R,\delta)$-truncated
  trajectory through $z_0$ is closed and admissible on $[0,2\pi]$.
\end{lemma}

\begin{proof}
  \uses{lem:mixed_sphere_margin_transport, lem:step_L1_reparam_relaxed,
        lem:step_error_expansion, thm:existence_of_zero}
  The same six-step assembly as the proof of the positive-stage
  \cref{lem:spherical_endpoint_winding} (Data $\to$ boundary margin $\to$
  reparametrization $\to$ winding $\to$ zero $\to$ admissibility), rebuilt
  from the sign-agnostic components rather than by invoking the positive
  wrapper (whose hypothesis \cref{def:sphere_four_vertex} we do not have).
  Two substitutions: the margins come from
  \cref{lem:mixed_sphere_margin_transport} (with $\kappa_0$ the mixed-sign
  lower bound of \cref{def:spherical_admissible_profile} instead of
  $\min\kappa>0$), and the $L^1$ reparametrization is
  \cref{lem:step_L1_reparam_relaxed} (continuity and periodicity of $\kappa$
  suffice; the step levels $a,b$ are chosen in the \emph{positive} part of
  the window, so the target step is the same positive-stage object and
  \cref{lem:step_error_expansion} applies unchanged). All other ingredients
  (\cref{def:spherical_endpoint}, \cref{lem:spherical_endpoint_continuousOn},
  \cref{lem:step_model_transport}, \cref{lem:winding_conj_loop},
  \cref{lem:winding_number_c_perturb}, \cref{thm:existence_of_zero}) are
  already sign-agnostic and are consumed verbatim.
\end{proof}

\begin{lemma}[mixed-sign simplicity]
  \label{lem:mixed_spherical_simplicity}
  \lean{Gluck.mixed_spherical_simplicity}
  \uses{lem:mixed_spherical_endpoint_winding, lem:reconstruction_ode,
        lem:spherical_trajectory_speed,
        lem:spherical_trajectory_eq_reconstruct,
        lem:is_simple_closed_const_add,
        cor:positive_curvature_implies_simple, def:error_vector}
  The closed admissible mixed-sign trajectory of
  \cref{lem:mixed_spherical_endpoint_winding} (periodically extended) is
  simple.
\end{lemma}

\begin{proof}
  \uses{lem:spherical_trajectory_eq_reconstruct,
        lem:is_simple_closed_const_add,
        cor:positive_curvature_implies_simple}
  The S2-E chain verbatim, through the relaxed (post-re-sign)
  \cref{lem:reconstruction_ode}: the exposed derivative conjunct gives
  $\tilde z'=\rho(t)e^{it}$ with trajectory speed
  $\rho(t)=q_{\kappa\circ h_1}(t,\tilde z(t))$ continuous, $2\pi$-periodic and
  $>0$ — positivity from the admissibility margin (denominator $\ge2\delta$),
  the sign of $\kappa$ irrelevant, exactly as in
  \cref{lem:spherical_trajectory_speed} whose proof never consumes $\kappa>0$.
  Then \cref{lem:spherical_trajectory_eq_reconstruct} identifies $\tilde z$
  with $\tilde z(0)+\mathrm{reconstruct}\,\rho$, closure kills the error
  vector, \cref{cor:positive_curvature_implies_simple} (public
  \texttt{isSimpleClosed\_reconstruct}) gives simplicity of
  $\mathrm{reconstruct}\,\rho$, and \cref{lem:is_simple_closed_const_add}
  transfers it to $\tilde z$. At stage-2 dispatch time the two trajectory
  lemmas should be stated (or generalized in place) with the relaxed
  hypothesis so the same Lean declarations serve both stages; if they land
  positive-only in S2-E, stage 2 re-signs them the same way as
  \cref{lem:reconstruction_ode}.
\end{proof}

\begin{theorem}[spherical converse, mixed sign]
  \label{thm:spherical_converse}
  \lean{Gluck.sphericalConverse}
  \uses{def:mixed_sign_sphere_four_vertex, def:realizes_spherical_curvature,
        lem:mixed_spherical_endpoint_winding, lem:mixed_spherical_simplicity,
        lem:reconstruction_ode, lem:realizes_spherical_comp,
        lem:spherical_circle_realizes}
  If $\kappa$ satisfies the mixed-sign spherical four-vertex hypothesis
  (\cref{def:mixed_sign_sphere_four_vertex}), then there is a simple closed
  curve $z$ confined to the open disk realizing $\kappa$ as its spherical
  geodesic curvature:
  \[
    \texttt{MixedSignSphereFourVertex }\kappa\;\Longrightarrow\;
    \exists\,z,\ \texttt{IsSimpleClosed }z\ \wedge\
    \texttt{RealizesSphericalCurvature }z\,\kappa .
  \]
\end{theorem}

\begin{proof}
  Same assembly as \cref{thm:spherical_converse_pos}: constant branch by the
  centered circle (\cref{lem:spherical_circle_realizes}; the constant is
  positive by hypothesis); non-constant branch by
  \cref{lem:mixed_spherical_endpoint_winding}, truncation removal by the
  relaxed-hypothesis variant of \cref{lem:reconstruction_ode} (speed
  positivity from the admissibility margin $\delta$, not from $\kappa>0$),
  pullback along $h_1^{-1}$ by \cref{lem:realizes_spherical_comp}, and
  simplicity by \cref{lem:mixed_spherical_simplicity}.
\end{proof}
